% !TEX root = ../../../main.tex
% Sources: LEC03, IMG_9133--IMG_9134; Week 3 topic summary.
\section{Polynomial functions and absolute value}
\label{sec:analysis-i:elementary-functions}

A \emph{polynomial function} has the form
\[
  p(x)=a_0+a_1x+\cdots+a_nx^n,\qquad a_0,\ldots,a_n\in\R.
\]
If $a_n\ne0$, its degree is $n$. It is defined for every $x\in\R$.
A \emph{rational function} is a quotient $p(x)/q(x)$ of polynomials,
with $q$ not the zero polynomial. Its domain excludes the zeros of $q$.
For example, the expression $(x^2-1)/(x-1)$ equals $x+1$ on its domain,
but the original quotient is undefined at $x=1$.

\subsection{Absolute value and distance}
The absolute value function, introduced before convergence in
\autoref{sec:analysis-i:convergence-preview}, is
\[
  |x|=\begin{cases}
    x,&x\geq0,\\
    -x,&x<0.
  \end{cases}
\]
It satisfies $|x|\geq0$, with equality exactly when $x=0$, and
$-|x|\leq x\leq|x|$. \autoref{fig:absolute-value} shows its graph.

\autoref{thm:analysis-i:absolute-value-inequalities} established
the product rule and triangle inequality before their use in limits.
Applying the triangle inequality to $x=(x-y)+y$, and then interchanging
$x$ and $y$, also gives the useful reverse triangle inequality
$\bigl||x|-|y|\bigr|\leq|x-y|$.

The following exercise asks for this argument and its equality case.
% !TEX root = ../../hw03.tex
% Source: HW3 Intro to Math Analysis I (1).pdf, page 2, question 7.
\begin{exercise}
  \label{ex:hw03:reverse-triangle-inequality}
  \solutionlink{hw03:reverse-triangle-inequality}
  Use the triangle inequality to prove
  \[
    \bigl||x|-|y|\bigr|\leq|x-y|,\qquad x,y\in\R.
  \]
  When will this become an equality?
\end{exercise}


\subsection{Positive and negative parts}
Define two nonnegative quantities
\[
  x^+=\max\{x,0\}=\frac{|x|+x}{2},\qquad
  x^-=\max\{-x,0\}=\frac{|x|-x}{2}.
\]
Checking separately $x\geq0$ and $x<0$ gives
\[
  x=x^+-x^-,\qquad |x|=x^++x^-,\qquad x^+x^-=0.
\]
The negative part records the \emph{magnitude} of the negative contribution,
so it is nonnegative. For a real-valued function $f$, apply these formulas
pointwise to obtain $f=f^+-f^-$ and $|f|=f^++f^-$.
\autoref{fig:positive-negative-parts} replaces the lecture's freehand
sketch with a specific example.
\begin{figure*}[htbp]
  \centering
  \begin{tikzpicture}[analysis diagram,x=21mm,y=21mm]
  \foreach \shift/\heading in {0/{$f(x)=x(1-x)$},2.2/{$f^+(x)$},4.4/{$f^-(x)$}} {
    \begin{scope}[shift={(\shift,0)}]
      \draw[->] (-0.5,0) -- (1.5,0) node[right] {$x$};
      \draw[->] (0,-0.8) -- (0,0.95) node[above] {$y$};
      \fill (0,0) circle[radius=1.6pt];
      \fill (1,0) circle[radius=1.6pt];
      \node[below left,analysis label] at (0,0) {$0$};
      \node[below left,analysis label] at (1,0) {$1$};
      \node[above] at (0.5,1.15) {\heading};
    \end{scope}
  }
  \draw[analysis curve,<->,domain=-0.45:1.45,samples=60]
    plot (\x,{\x*(1-\x)});
  \begin{scope}[shift={(2.2,0)}]
    \draw[analysis curve,<-] (-0.5,0) -- (0,0);
    \draw[analysis curve,domain=0:1,samples=40] plot (\x,{\x*(1-\x)});
    \draw[analysis curve,->] (1,0) -- (1.5,0);
  \end{scope}
  \begin{scope}[shift={(4.4,0)}]
    \draw[analysis curve,<-,domain=-0.45:0,samples=25] plot (\x,{\x*(\x-1)});
    \draw[analysis curve] (0,0) -- (1,0);
    \draw[analysis curve,->,domain=1:1.45,samples=25] plot (\x,{\x*(\x-1)});
  \end{scope}
\end{tikzpicture}

  \caption[Positive and negative parts of a function.]{Positive and negative parts of the illustrative function
    $f(x)=x(1-x)$. The part below the axis is reflected upward in $f^-$;
    in $f^+$ it is replaced by zero. All three panels use the same scale.}
  \label{fig:positive-negative-parts}
\end{figure*}
